\newpage
\phantomsection
\label{prf:quotient-rule}

\begin{remark*}[Return]
\hyperref[thm:quotient-rule]{Return to Theorem}
\end{remark*}

\begin{theorem*}[Quotient Rule]
Let $I \subseteq \mathbb{R}$ be an open interval, let $c \in I$, and let
$f, g : I \to \mathbb{R}$ be differentiable at $c$ with $g(c) \neq 0$. Then
$f/g$ is differentiable at $c$ and
\[
  \left(\frac{f}{g}\right)'(c) = \frac{f'(c)g(c) - f(c)g'(c)}{(g(c))^2}.
\]
\end{theorem*}

\begin{proof}
\textbf{Professional Standard Proof.}~\\
Set $h := f/g$. Since $g$ is differentiable at $c$, it is continuous at $c$, so $\lim_{x \to c} g(x) = g(c) \neq 0$. By the non-vanishing limit theorem, $g(x) \neq 0$ for all $x$ in some deleted neighbourhood of $c$, so $h$ is defined near $c$. For $x \neq c$ with $g(x) \neq 0$:
\[
  \frac{h(x) - h(c)}{x - c}
  = \frac{f(x)g(c) - f(c)g(x)}{g(x)g(c)(x - c)}
  = \frac{f(x)g(c) - f(c)g(c) + f(c)g(c) - f(c)g(x)}{g(x)g(c)(x - c)}.
\]
Factoring yields
\[
  \frac{1}{g(x)g(c)}
    \left[
      g(c) \cdot \frac{f(x) - f(c)}{x - c}
      - f(c) \cdot \frac{g(x) - g(c)}{x - c}
    \right].
\]
Taking the limit as $x \to c$ and applying limit laws:
\[
  h'(c) = \frac{1}{g(c)^2}[g(c) \cdot f'(c) - f(c) \cdot g'(c)] = \frac{f'(c)g(c) - f(c)g'(c)}{(g(c))^2}.
\]
\end{proof}

\begin{proof}
\textbf{Detailed Learning Proof.}~\\

\textbf{Step 1.} Form the difference quotient of $h = f/g$ and find a common denominator. Since $g$ is differentiable at $c$, it is continuous at $c$, so $\lim_{x \to c} g(x) = g(c) \neq 0$. By the non-vanishing limit theorem, $g(x) \neq 0$ for all $x$ in some deleted neighbourhood of $c$, ensuring $h$ is defined near $c$. For $x \neq c$ with $g(x) \neq 0$:
\[
  \frac{h(x) - h(c)}{x - c}
  = \frac{\frac{f(x)}{g(x)} - \frac{f(c)}{g(c)}}{x - c}
  = \frac{f(x)g(c) - f(c)g(x)}{g(x)g(c)(x - c)}.
\]

\textbf{Step 2.} Insert the auxiliary term $f(c)g(c)$ to decouple the increments of $f$ and $g$. Add and subtract $f(c)g(c)$ in the numerator:
\[
  \frac{f(x)g(c) - f(c)g(c) + f(c)g(c) - f(c)g(x)}{g(x)g(c)(x - c)}.
\]
Factor each pair and separate over the denominator:
\[
  = \frac{1}{g(x)g(c)}
    \left[
      g(c) \cdot \frac{f(x) - f(c)}{x - c}
      - f(c) \cdot \frac{g(x) - g(c)}{x - c}
    \right].
\]
This isolates the difference quotients of $f$ and $g$, which approach $f'(c)$ and $g'(c)$ respectively.

\textbf{Step 3.} Pass to the limit using limit laws. By the sum, product, and quotient rules for limits:
\[
  \lim_{x \to c} \frac{h(x) - h(c)}{x - c}
  = \frac{1}{g(c) \cdot g(c)}
    \left[
      g(c) \cdot f'(c) - f(c) \cdot g'(c)
    \right],
\]
where $\lim_{x \to c} g(x) = g(c)$ by continuity of $g$, and the difference quotient limits equal $f'(c)$ and $g'(c)$ by definition of differentiability. Therefore:
\[
  h'(c) = \frac{g(c)f'(c) - f(c)g'(c)}{(g(c))^2}.
\]
\end{proof}

\begin{remark*}[Proof structure]
The proof follows a direct approach mirroring the product rule strategy. After combining fractions in the difference quotient of $f/g$, the add-and-subtract technique with the auxiliary term $f(c)g(c)$ decouples the increments of $f$ and $g$ in the numerator, isolating the individual difference quotients. Continuity of $g$ at $c$ provides both the limit $\lim g(x) = g(c)$ in the denominator and the guarantee that $g(x) \neq 0$ near $c$ so the quotient is well-defined throughout.
\end{remark*}

\begin{remark*}[Dependencies]~\\
\begin{itemize}
  \item \hyperref[def:derivative-at-point]{Definition of Differentiability at a Point}
  \item \hyperref[thm:diff-implies-cont]{Differentiability Implies Continuity}
  \item \hyperref[thm:limit-nonvanishing]{Non-vanishing Limit}
  \item \hyperref[thm:limit-sum]{Sum Rule for Limits}
  \item \hyperref[thm:limit-product]{Product Rule for Limits}
  \item \hyperref[thm:limit-quotient]{Quotient Rule for Limits}
\end{itemize}
\end{remark*}